\secnumberlesssection{GLOSARIO}

\begin{description}
    \item[CPC] Componente común de viajeros (\textit{Common Part of Commuters}); medida de solapamiento entre flujos observados y predichos.
    \item[CRS] Sistema de referencia de coordenadas (\textit{Coordinate Reference System}).
    \item[DPA] División político-administrativa.
    \item[DTPM] Directorio de Transporte Público Metropolitano.
    \item[EOD] Encuesta Origen-Destino.
    \item[EPSG] Identificador del registro de sistemas de referencia geodésicos originado por el \textit{European Petroleum Survey Group}.
    \item[GAT] Red de atención en grafos (\textit{Graph Attention Network}).
    \item[H3] Sistema jerárquico de indexación geoespacial utilizado para definir unidades de origen y destino.
    \item[INE] Instituto Nacional de Estadísticas.
    \item[MAUP] Problema de la unidad espacial modificable (\textit{Modifiable Areal Unit Problem}); sensibilidad del análisis a la agregación y delimitación de unidades.
    \item[MLP] Perceptrón multicapa (\textit{Multilayer Perceptron}).
    \item[OD] Origen-destino.
    \item[OSM] OpenStreetMap.
    \item[POI] Punto de interés (\textit{Point of Interest}).
    \item[RMSE] Raíz del error cuadrático medio (\textit{Root Mean Squared Error}).
    \item[SECTRA] Secretaría de Planificación del Transporte.
    \item[SHAP] \textit{SHapley Additive exPlanations}; enfoque de atribución de predicciones basado en valores de Shapley.
    \item[UTM] Proyección universal transversal de Mercator (\textit{Universal Transverse Mercator}).
    \item[WGS 84] Sistema Geodésico Mundial de 1984 (\textit{World Geodetic System 1984}).
    \item[XAI] Inteligencia artificial explicable (\textit{eXplainable Artificial Intelligence}).
\end{description}
